\section{From complement failure to the pinning input}\label{sec:H}

We prove Theorem~\ref{thm:birational} by translating a failure of bounded
klt complements into the sequence excluded by Theorem~\ref{thm:G-pinning}.
The construction uses a canonical cover when the birational morphism is an
isomorphism, and a signed anti-canonical torsor otherwise.  We then apply
Chen's equivalence and justify the passage between algebraically closed
fields of characteristic zero.

\subsection{A prime saturating the complete systems}

Fix a dimension $d>0$ and $0<\epsilon\le1$, and suppose the birational
complement assertion fails over $\C$.  Birkar's relative complement
theorem gives an integer $b=b(d)>0$ such that every example has an lc
$b$-complement near its specified point \cite[Theorem~1.8]{BirkarComplements}.
The hypotheses are precisely Fano type, the zero boundary, and nef
$-K_X$; in particular no generalized nef-part index is being imposed.
Choose a sequence with no klt $M_i$-complement, where $M_i>1$ tends to
infinity, $b\mid M_i$, and every fixed positive integer divides $M_i$
eventually.  For example, sufficiently large factorials have these
properties.  Failure of a uniform index allows this choice for each
prescribed $M_i$.

We may assume $-K_X$ is ample over $Z$.  On a variety of Fano type, the
relative Mori dream space consequences of the klt MMP imply that nef
$-K_X$ is semiample \cite{BCHM}.  Since $f$ is birational, this divisor
is relatively big.  Its ample model is thus a birational contraction
$g:X\to X^{\mathrm a}$ over $Z$.  A divisible multiple of $-K_X$ is the
pullback of its ample model divisor.  Pushing compatible canonical
divisors forward identifies the latter with the corresponding multiple
of $-K_{X^{\mathrm a}}$, so $K_X=g^*K_{X^{\mathrm a}}$ after compatible
choices.  Equivalently, a possible rational principal adjustment descends
because the function fields agree.  Thus the contraction is crepant,
$X^{\mathrm a}$ remains $\epsilon$-lc and of Fano type over $Z$, and
klt complements on it pull back to klt complements on $X$
\cite{ProkhorovShokurov}.  Applying the lc complement theorem on this
model and replacing $X$ by it preserves the failure under consideration.

Shrink $Z$ to an affine neighbourhood of the specified point.  For every
integer $k$, a section of $\mathcal O_X(-kK_X)$ is represented by a
rational function $s$ with
\begin{equation}\label{eq:H-section-divisor}
 \Delta_s=\divisor(s)-kK_X\ge0.
\end{equation}
The degree-$b$ complement is represented by $s_b$ and satisfies
$(X,\Delta_{s_b}/b)$ lc.  Here and below all section systems are complete
relative systems on this neighbourhood.

\begin{proposition}[Factorial saturation]\label{prop:H-saturation}
For each example there is a primitive divisorial order $w$ over $X$ whose
image centre on $Z$ contains the specified point and for which
\begin{equation}\label{eq:H-saturation}
 w(\Delta_s)\ge k A_X(w)\quad
       (k>0,\ k\mid M_i),\qquad
 w(\Delta_{s_b})=bA_X(w).
\end{equation}
After replacing the example by a geometric generic slice, one may assume
that the image centre is a closed point $z$, that
$1\le\dim X=\dim Z\le d$, and that all the inequalities hold for the
complete local systems at $z$.
\end{proposition}
\begin{proof}
The coherent sheaf $f_*\mathcal O_X(-M_iK_X)$ has finitely many
generators on the affine base.  Resolve their fractional ideal of ratios
to a fixed section, together with the relevant divisors.  On the resulting
smooth model $p:Y\to X$, their pulled-back divisor system has a common
fixed $\Q$-divisor $F$ and a basepoint-free residual system.  This
construction is valid even if $M_iK_X$ is not Cartier: each section divisor
is $\Q$-Cartier, their differences are principal, and principalizing the
fractional ratio ideal makes the moving differences Cartier.  The fixed
divisor has order equal to the minimum of the section orders at every
prime on the model.

The section $s_b^{M_i/b}$ belongs to this system.  Therefore
$F/M_i\le p^*(\Delta_{s_b}/b)$, and the fixed-order discrepancy
\[
 A_X(E)-\frac{\ord_E(F)}{M_i}
\]
is nonnegative for every prime $E$ over $X$.  Take the support on $Y$ to
be simple normal crossings.  A general residual member, with coefficient
$1/M_i<1$, is transverse to this support.  If every fixed-order
discrepancy were strictly positive over a neighbourhood of the specified
point, this general member would give a klt $M_i$-complement there.
The only possible zero discrepancies of the fixed log smooth data come
from strata containing a coefficient-one component.  There are finitely
many such components, and their images in $Z$ are closed because
$Y\to Z$ is proper.  Thus the absence of a klt complement on every
neighbourhood implies that one component $E$ with image containing the
specified point has
$\ord_E(F)=M_i A_X(E)$.  Set $w=\ord_E$.

For $k\mid M_i$ and any section $s$ in degree $k$, its power
$s^{M_i/k}$ is in degree $M_i$.  Hence
$(M_i/k)w(\Delta_s)\ge\ord_E(F)=M_i A_X(w)$.
The same inequality for $s_b$, together with lc of the chosen complement,
gives equality in degree $b$.  These statements also apply to local
sections: coefficients in the local base ring have nonnegative $w$-order,
and shrinking the base only inverts functions of order zero whenever
the centre still contains the specified point.

Write $W$ for the image centre of $w$ on $Z$.  It is a proper closed
subset: $w(\Delta_{s_b})>0$ and a divisorial centre cannot dominate the
birational base.  If $q=\dim W>0$, choose a general linear projection
$Z\to\mathbf A^q$ that is dominant and generically finite on $W$.
Pass to the geometric generic fibre and follow a component of $E$.
The image of that component in the new base is a closed point, and the
new dimension is $\dim X-q\ge1$.

Here are the required preservation details.  The prime $E$ is horizontal
over $\mathbf A^q$.  After shrinking its base, the chosen resolution and
its finitely many relevant strata are smooth wherever they dominate this
base.  Dividing a top form by a base volume form therefore restricts
compatible canonical divisors to the generic fibre, preserving the
log discrepancies and section-divisor orders.  Extending the generic
constant field separably does not ramify these prime orders: the
nonzero discriminant of a finite separable constant extension is a unit
in their valuation rings.  Geometric integrality can be obtained by
choosing the corresponding components after this extension.

Divisorial sheaves restrict to the corresponding divisorial sheaves by
localization and separable scalar extension, as is seen on the smooth
locus and then in codimension one.  Flat base change for the proper morphism
shows that the original generators generate the new complete systems;
only the finitely many degrees dividing $M_i$ and the degree $b$ are
needed here.  Their inequalities hence hold for every local section at
the new image point.  Normality, birationality, the discrepancy bound,
relative ampleness of $-K_X$, and a Fano type witness are preserved by
this operation.  Thus all the conclusions survive the slice.  The
algebraically closed generic field of the complex base is isomorphic
to $\C$; alternatively the finite data can again be embedded into $\C$.
We retain the notation $X,Z,z,w$ for these new data.
\end{proof}

After a subsequence the new dimension and which of the following two
constructions applies are fixed.  We keep the original bounded integer
$b$, even if the dimension has decreased.

\subsection{The canonical cover and the signed torsor}

\begin{lemma}[The isomorphic case]\label{lem:H-canonical-cover}
Suppose $f$ is an isomorphism near $z$.  There is a normal integral
cyclic cover $S_0\to X$, with a unique point $o$ over $z$, a
quasi-\'{e}tale cyclic action fixing $o$, and a divisor-free canonical
generator of character one.  It has a regular eigenfunction
$f_b=s_bt^b$ of character $b$ and a divisorial valuation $w_0$ centred
at $o$ such that, with $A_0=A_{S_0}(w_0)>0$,
\begin{equation}\label{eq:H-upstairs-saturation}
 (S_0,\divisor(f_b)/b)\text{ is lc},\qquad
 w_0(f_b)=bA_0,\qquad
 w_0(a)\ge kA_0
\end{equation}
for every regular eigenfunction $a$ of character $k$ when
$k>0$ and $k\mid M_i$.  The variety $S_0$ is klt.
\end{lemma}
\begin{proof}
Let $r$ be the local Cartier index of $K_X$ and shrink so that
$rK_X=\divisor(c)$.  Normalize $X$ in
\[
 t^r=c^{-1},
\]
with the cyclic group acting on $t$ by its standard character.  The cover
is integral of degree $r$.  Indeed a smaller Kummer degree would make
$c^{-1}$ a $p$th power for a prime $p\mid r$, and then
$(r/p)K_X$ would be principal, contrary to minimality.  At a prime of
$X$, $K_X$ is Cartier and the exponent of $c$ is divisible by $r$.
The codimension-one Kummer calculation shows that the cover is
quasi-\'{e}tale.  If $\divisor(\omega_X)=K_X$, the form
$t\omega_X$ has divisor zero upstairs.  Finite crepant pullback proves
that $S_0$ is klt.

A fibre of this finite quotient is a single orbit: invariant functions
separate distinct finite orbits.  If the fibre over $z$ had more than one
point, regular functions separating its points would have a nontrivial
eigencomponent nonzero at one, hence at every, point of that orbit.
After shrinking by its invariant norm, it would be an eigenunit.  Write
it $s t^j$ with $0<j<r$.  The codimension-one calculation gives
$\divisor(s)=jK_X$, contradicting the minimal index.  Thus the fibre is
one fixed point.

For every integer $k$, the regular functions of character $k$ are
precisely $s t^k$ with $s\in\mathcal O_X(-kK_X)$; changing $k$ by
$r$ is accommodated by $t^r=c^{-1}$.  Their divisors are the pullbacks
of $\Delta_s$.  Extend $w$ to a primitive divisorial valuation $w_0$
on the cover.  Its centre is $o$.  If its restriction is $e w$, finite
crepant pullback gives $A_{S_0}(w_0)=eA_X(w)$ and
$w_0(s t^k)=e w(\Delta_s)$.  Equation~\eqref{eq:H-saturation} proves
all the assertions.
\end{proof}

For the second construction, retain the grading variable $t$ as an
indeterminate, and write $K=\C(X)=\C(Z)$.
Choose a rational top form $\omega_X$ with $\divisor(\omega_X)=K_X$.
The positive degrees record the anti-canonical section inequalities.
We also keep negative degrees: after restriction to a finite scalar
group, one character space contains degrees of both signs.  Bounds on
both signs will therefore be needed to control every section of a
fixed character.

\begin{lemma}[The signed anti-canonical torsor]\label{lem:H-signed-torsor}
Suppose $f$ is not an isomorphism near $z$.  The signed algebra
\begin{equation}\label{eq:H-signed-algebra}
 R=\bigoplus_{k\in\Z}
       H^0(X,\mathcal O_X(-kK_X))t^k,\qquad S_0=\Spec R,
\end{equation}
over the affine base is a finitely generated normal domain of dimension
$\dim X+1$. Write $R_k=H^0(X,\mathcal O_X(-kK_X))t^k$ for its
homogeneous part of degree $k$. Its \emph{positive chamber} is the open
\[
 S_0^+=\bigcup_{k>0}\ \bigcup_{0\ne a\in R_k}D_{S_0}(a),
\]
where $D_{S_0}(a)$ is the principal open on which $a$ does not vanish.
This chamber has complement of codimension at least two. The form
\begin{equation}\label{eq:H-torsor-form}
 \Omega_0=t\omega_X\wedge d\log t
\end{equation}
is a divisor-free canonical generator of grading weight one, and every
finite scalar grading subgroup acts quasi-\'{e}talely.

The variety $S_0$ is klt, and
$(S_0,\divisor(s_bt^b)/b)$ is lc.  There is a grading-fixed closed point
$o$ over $z$ and a rational Gauss valuation $w_0$ centred there which
satisfies \eqref{eq:H-upstairs-saturation} for every actual positive
degree $k\mid M_i$.  Moreover there are numbers $\alpha>0$ and
$F_w>\alpha$, depending on the example, for which every nonzero
homogeneous section satisfies
\begin{equation}\label{eq:H-two-slopes}
 w_0(s t^k)\ge
 \begin{cases}
 k\alpha,&k>0,\\
 |k|(F_w-\alpha),&k<0.
 \end{cases}
\end{equation}
\end{lemma}
\begin{proof}
\textit{Finite generation and normality.}
The positive algebra is finitely generated by relative ampleness of
$-K_X$.  For the negative algebra, take a small $\Q$-factorialization
and a $K$-MMP over $Z$.  Fano type gives termination and a semiample
output, since every divisor is relatively big for a birational
contraction \cite{BCHM}.  The integral-degree section spaces are
preserved in these steps.  To check this without assuming a Cartier
index in degree one, on a common model write
$p^*K_{X_1}=q^*K_{X_2}+E$ for a $K$-MMP step, where $E$ is effective
and $q$-exceptional.  An effective section divisor
$\divisor(s)+mK_{X_1}$ pushes to an effective section divisor on $X_2$.
Conversely, the pullback of an effective section divisor on $X_2$,
plus $mE$, is effective, and pushes to its section divisor on $X_1$.
These are rational pullbacks, valid for every integral $m>0$ because
the divisors involved are $\Q$-Cartier; one can equivalently verify
the inequalities after taking a Cartier power.  The initial small
modification identifies the divisorial sheaves in codimension one.
On the semiample output $Y$, let $g:Y\to Y^{\mathrm c}$ be the
semiample contraction and choose $rK_Y=g^*L$ with $L$ relatively ample.
For each $0\le j<r$, the projection formula identifies the residue
module with
\[
 \bigoplus_{a\ge0}
 H^0\bigl(Y^{\mathrm c},g_*\mathcal O_Y(jK_Y)\otimes
                           \mathcal O_{Y^{\mathrm c}}(aL)\bigr).
\]
The pushforward is coherent by properness, so this is finite over the
ample section ring.  The Cartier Veronese is finitely generated and
the finitely many residue modules are finite over it.  The negative
algebra is therefore finitely generated.  The union of generators of
the two half-algebras generates $R$.

The ring $R$ is the intersection of $K[t,t^{-1}]$ with the Gauss
valuation rings specified, for every prime $D$ of $X$, by
\[
 \ord_D(s)-k\operatorname{coeff}_D(K_X)\ge0
 \quad\text{on }s t^k.
\]
This follows from the codimension-one description of a divisorial
sheaf, and proves normality.  Its fraction field is $K(t)$: the
fraction field of degree zero is $K$, and the generic fibre contains
every integral degree.  In particular $\dim S_0=\dim X+1$.

\medskip
\noindent\textit{The positive chamber and the canonical form.}
We prove that every prime divisor meets the positive chamber, where
the canonical form can be computed on a Laurent bundle.  Suppose
that a prime divisor $P$ lies outside it.  The connected grading torus
preserves every irreducible component of this complement, so $P$ is
homogeneous.  It cannot contain all nonzero graded pieces.  Indeed the
locus where they all vanish is a closed subscheme of $Z$ missing its
generic point, where the signed algebra is $K[t,t^{-1}]$.  That locus
has dimension at most $\dim Z-1=\dim S_0-2$.

Choose a negative homogeneous element $q$ of degree $-m$ not in $P$.
Set
\[
 A=R[q^{-1}]=R_{\le0}[q^{-1}],\qquad C=A_0.
\]
The equality follows by multiplying any positive-degree element by
sufficiently many powers of $q$.  The ring
$C[q,q^{-1}]$ is a Laurent polynomial ring: its nonzero powers have
different degrees.  For a homogeneous $a\in A_k$ the identity
\begin{equation}\label{eq:H-laurent-integrality}
 a^m q^k\in C,\qquad
 a^m=(a^m q^k)q^{-k}
\end{equation}
proves integrality over this Laurent ring.  Both formulas are valid for
negative $k$, because $q$ is a unit on this chart.  Finite generation
then makes $A$ finite over $C[q,q^{-1}]$.
The ring $C$ contains the affine base ring $R_0$ and lies in $K$, so
$\Frac C=K$.

The height-one homogeneous prime induced by $P$ contracts to a
height-one homogeneous prime of $C[q,q^{-1}]$.  Since $q$ is invertible,
that contracted prime is generated by its intersection with $C$ and
lies over a prime divisor of $\Spec C$, a chart of the negative Proj.
This also excludes trivial restriction to $K$: there is no homogeneous
prime divisor solely in the invertible Laurent direction.  The negative
Proj is the relatively ample $K$-model obtained by the small modification,
the $K$-MMP and its semiample contraction.  None of these operations
extracts divisors.  The prime in question therefore has a strict
transform $D$ on $X$; its restriction has the form
$w'=e\ord_D$, $e>0$.

There are sections of both sufficiently divisible signs whose
$\Delta$-orders at $D$ are zero.  For the positive sign this is relative
generation of a multiple of $-K_X$.  For the negative sign use a
generated multiple of $K$ on its ample model avoiding the surviving
prime.  The birational map identifies the two generic DVRs and their
compatible canonical coefficients.  Its section lifts through the
preserved section space, and the effective MMP correction has zero
coefficient at this strict divisor.  Thus its $\Delta$-order on $X$ is
also zero.  Relative generation supplies both sections globally over
the affine base with the indicated degree signs.

Let $v=\ord_P$, and put
$\beta=v(t)+w'(K_X)$, where $w'(K_X)$ means rational pullback order.
For a homogeneous element,
\begin{equation}\label{eq:H-gauss-order}
 v(s t^k)=w'(\Delta_s)+k\beta.
\end{equation}
Nonnegativity of $v$ on the two zero-multiplicity sections forces
$\beta\ge0$ and $\beta\le0$.  Hence $\beta=0$, and the positive
section just chosen has value zero, contrary to $P$ being outside the
positive chamber.  This proves the codimension assertion.

On the positive chamber the signed algebra is the usual Laurent
divisorial bundle over $X$.  To see this on a Proj chart, localize at a
divisible positive section.  Poles along its complementary Cartier
divisor can be cleared by its powers, identifying every degree piece
with the corresponding localized divisorial sheaf on $X$.  A Cartier
Veronese is a line bundle with its zero section removed, and the full
chart is finite over that bundle by taking powers.  In particular the
inverse image of a codimension-at-least-two subset of $X$ has codimension
at least two.  At a regular codimension-one point write
$K_X=\divisor(c)$; then $ct$ is an invertible Laurent coordinate and
\[
 t\omega_X\wedge d\log t
   =(ct)(\omega_X/c)\wedge d\log(ct).
\]
This is a nowhere-vanishing top form, and every nonidentity finite
scalar acts freely there.  The established codimension-two complements
extend these conclusions to the assertions about $\Omega_0$ and
quasi-\'{e}taleness on $S_0$.

\medskip
\noindent\textit{Discrepancies and the lc boundary.}
The form also controls valuations whose centres lie outside the positive
chamber.  We will use its discrepancy estimate both for the boundary
defined by $s_bt^b$ and for the finite quotients in
Lemma~\ref{lem:H-padding}.  If a divisorial multiple $v$ centred on $S_0$ has
nontrivial restriction $w'$ to $K$, the relative Abhyankar inequality
makes $w'$ a divisorial multiple: its value group is discrete, and
adjoining $t$ increases residue transcendence degree by at most one,
forcing its residue field to have transcendence degree $\dim X-1$ over
$\C$.
Its centre exists on $X$ by properness.
Relative generation of a positive Cartier multiple gives
$\beta=v(t)+w'(K_X)\ge0$.  The top form gives
\begin{equation}\label{eq:H-torsor-discrepancy}
 A_{S_0}(v)\ge A_X(w')+\beta,
\end{equation}
with equality for a Gauss extension.  One can compute this on a smooth
model at the prime defining $w'$.  Choose a uniformizer $x$ and
separating residue coordinates $y_1,\ldots,y_{\dim X-1}$.  The ratio of
$\omega_X$ to $d\log x\wedge dy_1\wedge\cdots\wedge dy_{\dim X-1}$
has $w'$-order $A_X(w')+w'(K_X)$.  After wedging with $d\log t$, the remaining
logarithmic differential has nonnegative log order at $v$.  Indeed, at
an upstairs uniformizer $\rho$, every logarithmic differential has at
most a simple $d\rho/\rho$ pole, at most one such factor survives a
wedge, and the differentials of the residue-coordinate units are
regular.  At a Gauss
valuation its residues are independent and this log order is zero.
Multiplication by $t$ yields \eqref{eq:H-torsor-discrepancy}.

If $w'$ is trivial, work on the smooth generic fibre
$\Spec K[t,t^{-1}]$; a centred valuation there has $v(t)=0$ and positive
discrepancy.  Thus $S_0$ is klt.  For a nontrivial restriction,
\eqref{eq:H-gauss-order}, \eqref{eq:H-torsor-discrepancy}, and lc of the
actual degree-$b$ complement give
\[
 A_{S_0}(v)-\frac{v(s_bt^b)}b
 \ge A_X(w')-\frac{w'(\Delta_{s_b})}b\ge0.
\]
For trivial restriction, $v(s_bt^b)=0$.  Consequently the displayed
boundary pair is lc.  It has the lc place constructed next and is not
being asserted to be klt.

\medskip
\noindent\textit{A centred order with two slopes.}
It remains to construct the closed point and its saturated valuation.
The positive slope will come from a choice of Gauss weight; the negative
slope comes from vanishing on the fibre.
Every irreducible component of the reduced fibre over $z$ has positive
dimension.  An isolated zero-dimensional component would disconnect
the fibre unless the whole fibre were finite; a proper birational
morphism to a normal base with a finite fibre is an isomorphism after
shrinking, the excluded case.  A negative-degree section, raised to a
Cartier power, restricts on each such projective component to a section
of an antiample line bundle.  It is therefore zero on every component.
Its effective divisor contains the entire reduced fibre, so its order
at the specified $w$ is strictly positive.  This argument only concerns
reduced supports and also covers nonreduced fibres.

Choose finitely many negative-ring generators $s_jt^{-m_j}$ over the
degree-zero ring and set
\[
 F_w=\min_j\frac{w(\Delta_{s_j})}{m_j}>0.
\]
Every negative homogeneous section is a homogeneous polynomial in these
generators with base coefficients of nonnegative $w$-order.  It follows
that $w(\Delta_s)\ge |k|F_w$ for $k<0$.
Choose a rational $0<\alpha<F_w$, and let $w_0$ be the Gauss extension
of $w$ with $w_0(t)+w(K_X)=\alpha$.  Formula
\eqref{eq:H-gauss-order} defines it, and gives
\eqref{eq:H-two-slopes}; for the positive sign use effectivity of
$\Delta_s$.  All nonzero degrees have strictly positive value, while
the degree-zero maximal ideal is $\mathfrak m_z$.  Thus $w_0$ is centred
at the fixed closed point with ideal generated by these elements.  A
rational Gauss extension of a divisorial valuation is itself divisorial
up to scaling.  Equality in \eqref{eq:H-torsor-discrepancy} gives
$A_0=A_X(w)+\alpha$, and \eqref{eq:H-saturation} now yields
\eqref{eq:H-upstairs-saturation} in actual positive degrees.
\end{proof}

\subsection{The quotient gap and matching the canonical character}

In the canonical-cover case retain its cyclic group.  In the torsor
case choose any finite scalar subgroup $\mu_r$; its order will be fixed
below, after the individual slopes in \eqref{eq:H-two-slopes} are known.
All characters in the next lemma are expressed relative to the standard
character of this cyclic group.

\begin{lemma}[Padding and the actual quotient discrepancy]\label{lem:H-padding}
In either construction, set
\[
 S=S_0\times\mathbf A^{b-1},
 \qquad (o,0)\in S,
\]
and give every added coordinate $u_j$ character one.  Then $S$ is klt,
the finite quotient $Y=S/\mu_r$ is $\epsilon$-lc, and the canonical
generator on $S$ has character $b$.  There is a regular function $h$ of
character $b$ such that $(S,\divisor(h))$ is lc and a centred
quasi-monomial valuation $T_0$ such that
\begin{equation}\label{eq:H-padding-equality}
 A_S(T_0)=T_0(h)=bA_0.
\end{equation}
\end{lemma}
\begin{proof}
The product is klt.  The action remains quasi-\'{e}tale: the fixed locus
of every nonidentity element has codimension at least two already in
$S_0$.  Wedging its canonical generator with
$du_1\wedge\cdots\wedge du_{b-1}$ gives a divisor-free generator of
character $b$.  A suitable power descends through the finite quotient,
so $K_Y$ is $\Q$-Cartier and the quotient map is crepant.

Quasi-\'{e}taleness alone would not preserve a prescribed discrepancy
gap.  Let $E$ be an actual prime divisor over $Y$, and normalize an
extension $v$ to $\C(S)$ so that its restriction to $\C(Y)$ is
$\ord_E$.  Finite crepant pullback, with homogeneous discrepancies,
gives $A_Y(E)=A_S(v)$.  Although $v$ may take fractional values upstairs,
the field $K=\C(X)$ is pointwise invariant and is contained in
$\C(Y)$.  If $w'=v|_K$ is nontrivial, its values are integral.  The
relative Abhyankar inequality then gives
\[
 w'=e\ord_F\quad\text{for a prime }F\text{ over }X
 \quad\text{and an integer }e\ge1.
\]
For the torsor, the log-form calculation above, with the additional
coordinate forms, yields
\[
 A_S(v)\ge A_X(w')+\beta+\sum_jv(u_j)
          \ge e\epsilon\ge\epsilon.
\]
Here $\beta\ge0$ by positive generation, each $u_j$ is regular, and
the non-Gauss logarithmic differential defect is nonnegative.  For the
canonical cover, first restrict to $\C(S_0)$ and use the product
log-form estimate; finite crepant pullback from $S_0$ to $X$ then gives
the same lower bound $A_X(w')$.

If $w'$ is trivial, the valuation lies over the common generic point
of the birational varieties $X$ and $Z$.  The generic quotient is smooth: the canonical cover there is
a cyclic torsor, and the signed construction is a Laurent line with a
free scalar action; their products with the coordinate representation
have smooth quotients.  Generic restriction of the ordinary prime
$E$ is still a primitive prime, since localization at the trivially
valued base field does not change its value group or uniformizer.
Its discrepancy on this smooth generic quotient is at least one, hence
at least $\epsilon$.  This proves the required gap for every actual
quotient prime.

Put $f_b=s_bt^b$ and take a general combination
\begin{equation}\label{eq:H-h-function}
 h=a_0 f_b+a_1u_1^b+\cdots+a_{b-1}u_{b-1}^b.
\end{equation}
We check lc using the ideal $I=(f_b,u_1^b,\ldots,u_{b-1}^b)$.
For any divisorial valuation $v$ with nontrivial restriction $v_0$ to
$\C(S_0)$, that restriction is divisorial up to scaling and the
product differential calculation gives
\[
 A_S(v)\ge A_{S_0}(v_0)+\sum_jv(u_j)
 \ge \frac{v(f_b)}b+\sum_jv(u_j)
 \ge \min\{v(f_b),bv(u_1),\ldots,bv(u_{b-1})\}.
\]
If the restriction is trivial, $v(f_b)=0$ and the same ideal inequality
holds because $S$ is klt.  Therefore $(S,I)$ is lc.  On a resolution of
the ideal, a general residual member is smooth and transverse to the
fixed log smooth boundary; with coefficient one this proves
$(S,\divisor(h))$ lc.  All summands of $h$ have character $b$.

Finally extend $w_0$ by the joint Gauss valuation giving every $u_j$
weight $A_0$.  It is centred at $(o,0)$ and quasi-monomial, and its
product discrepancy is $A_0+(b-1)A_0=bA_0$.  Every summand of $h$ has
that value.  Their initials do not cancel for a general choice of the
$a_j$ (the added coordinates are independent), giving
\eqref{eq:H-padding-equality}.  When $b=1$ there are no added
coordinates and $h$ is simply a nonzero scalar multiple of $f_b$.
\end{proof}

\begin{proposition}[The pinning input]\label{prop:H-g-input}
The failure sequence produces, after a subsequence, a fixed-dimensional
sequence satisfying every hypothesis of Theorem~\ref{thm:G-pinning}.
\end{proposition}
\begin{proof}
Let $\chi$ be character $b$ of the chosen cyclic group; it need not be
faithful.  Normalize $T=T_0/(bA_0)$.  The preceding lemma gives
\[
 A_S(T)=T(h)=1,
 \qquad (S,\divisor(h))\text{ lc},
 \qquad S/\mu_r\text{ is }\epsilon\text{-lc},
\]
and $h$ and the divisor-free canonical generator both have character
$\chi$.  It remains to check, for each fixed positive integer $m$, that
eventually
\begin{equation}\label{eq:H-final-saturation}
 T(a)\ge m
 \quad\text{for every regular eigenfunction }a\text{ of character }\chi^m.
\end{equation}

Expand a polynomial eigenfunction in monomials in the added coordinates.
A monomial of total degree $j$ has $T_0$-value $jA_0$.  If $j\ge bm$,
this already suffices, since its coefficient is regular and has
nonnegative value.  Otherwise its coefficient on $S_0$ has character
\[
 k=bm-j>0.
\]
In the canonical-cover case, that entire eigenspace is the degree-$k$
section space described in Lemma~\ref{lem:H-canonical-cover}.  Since
every fixed $k$ divides $M_i$ eventually, its coefficients have value
at least $kA_0$.  The monomial therefore has value at least
$(j+k)A_0=bmA_0$.

In the signed torsor, coefficients decompose into actual degrees
$k+\ell r$, $\ell\in\Z$.  For $\ell=0$ use the same factorial
saturation.  For $\ell\ne0$ use both strict slopes in
\eqref{eq:H-two-slopes}.  At example $i$, choose $r_i>i$ so large that
\begin{equation}\label{eq:H-cyclic-order}
 (r_i-i)\min\{\alpha,F_w-\alpha\}\ge iA_0.
\end{equation}
For $1\le k\le i$, every nonzero residue shift has
$|k+\ell r_i|\ge r_i-i$, so its value is at least $iA_0\ge kA_0$.
This also gives the desired monomial bound.  The choice is made after
$\alpha,F_w,A_0$ are known; no uniform lower bound on either slope is
needed.  For fixed $m$, eventually $bm\le i$, and every required
unshifted degree divides $M_i$.

A finite sum cannot have value smaller than the minimum of its terms.
The argument therefore proves \eqref{eq:H-final-saturation} in the
affine coordinate ring.  It also proves it in the local ring at the
fixed point.  Indeed any local function can be written with an invariant
unit denominator: multiply an initial denominator by its other group
translates and use its invariant norm.  For an eigenfunction the resulting
numerator has the same character.  The invariant denominator has value
zero, so the inequality is unchanged.

The dimensions satisfy
\[
 \dim S=
 \begin{cases}
 \dim X+b-1,&\text{canonical-cover case},\\
 \dim X+b,&\text{signed-torsor case},
 \end{cases}
 \qquad\text{and hence}\qquad \dim S\le d+b.
\]
They are thus bounded independently of $i$, and a subsequence has fixed
dimension.  The valuations are centred and quasi-monomial, the covers
are finite cyclic and quasi-\'{e}tale, and every remaining requirement
of Theorem~\ref{thm:G-pinning} was checked above.
\end{proof}

\subsection{Chen's reduction and change of field}

We use the following form of Chen's reduction.
For every positive integer $\ell$, Chen's Theorem~1.6, assertions (2) and (3),
identifies the following two assertions \cite{Chen2025}:
\begin{enumerate}[label=(\roman*)]
\item bounded monotonic klt complements for birational Fano type
contractions with zero boundary and $\dim Z\le\ell$;
\item the Cartier-section assertion for Fano type contractions with
$\dim Z\le\ell$, with a constant depending on the total-space dimension
and the discrepancy gap.
\end{enumerate}
In (i), the total-space and base dimensions are equal.  Consequently,
Theorem~\ref{thm:birational} in each of the finitely many positive
dimensions at most $\ell$, for every positive discrepancy parameter,
supplies (i); its dimension-zero case is
trivial.  One may take a common multiple of the finitely many complement
indices for each fixed discrepancy parameter.  Thus the birational
inputs have dimension at most $\ell$, whereas assertion (ii) also allows
total spaces of higher dimension.

The real-pair conventions in Chen's statement cause no extra hypothesis
in (i).  If an integral Weil divisor $K_X$ is $\R$-Cartier, express it
locally as a finite real combination of Cartier divisors.  Comparing Weil
coefficients gives a finite linear system with rational coefficients and
rational right-hand side.  A real solution implies a rational solution,
so $K_X$ is $\Q$-Cartier.  Similarly, a real effective log Fano witness
can be replaced by a rational one.  Keep its finite boundary support
and its zero coefficients.  On a finite open cover, choose Cartier
representations of its log canonical class, using the same boundary
coefficient variables on every chart.  These impose finitely many
rational affine equations.  Rational points are dense in their real
solution space.  Every rational solution defines a global rational
divisor whose log canonical divisor is locally $\Q$-Cartier; clearing
denominators on the finite cover makes it $\Q$-Cartier on $X$.  A rational
affine basis of the solution space therefore places all its log canonical
divisors in a fixed finite-dimensional real span of global $\Q$-Cartier
divisors.  Relative ampleness is open in that span.  On a fixed log
resolution, the klt condition and positivity of the remaining boundary
coefficients also persist under a sufficiently small perturbation.
Hence the rational birational
statement proves the full zero-boundary assertion needed in Chen's
real-pair formulation.

\begin{lemma}[Point ideals and change of field]\label{lem:H-field-reduction}
It suffices to prove Theorem~\ref{thm:main} over $\C$.  More precisely, a
uniform positive lower bound for the log canonical threshold of
$\mathfrak m_z\mathcal O_X$ implies the Cartier-section assertion after
replacing the bound by a smaller number less than one.  Both the hypotheses
and this ideal threshold are preserved by algebraically closed field
extensions of characteristic zero.
\end{lemma}
\begin{proof}
Shrink $Z$ to an affine neighbourhood of $z$, choose generators
$g_1,\ldots,g_N$ of $\mathfrak m_z$, and resolve both $(X,B)$ and the
pulled-back ideal.  If $p:Y\to X$ is the resulting log resolution, write
\[
 K_Y+B_Y=p^*(K_X+B),\qquad
 \mathfrak m_z\mathcal O_Y=\mathcal O_Y(-F).
\]
The ideal threshold is the minimum of
$A_{X,B}(E)/\ord_E(F)$ over the finitely many components of $F$ having
positive order.  Its interpretation is lc on the inverse image of a
neighbourhood of $z$.  Indeed the ideal is the unit ideal away from
$f^{-1}(z)$, and the original pair is already lc there.

For a general linear combination $g=\sum a_jg_j$, the pullback of
$\divisor(g)$ is $F+H$, where $H$ is a general member of a basepoint-free
system.  On a further restriction of the base, $H$ is smooth and transverse
to the fixed simple normal crossings support.  If $0<t<1$ is at most the
ideal threshold, then $(Y,B_Y+tF+tH)$ is lc, hence so is
$(X,B+t f^*\divisor(g))$.  Choose $g\ne0$.  Since $g\in\mathfrak m_z$,
$D=\divisor(g)$ is a nonzero effective Cartier divisor containing $z$.
Conversely, a Cartier divisor through $z$ has a local equation
$g\in\mathfrak m_z$, so that $v(\mathfrak m_z\mathcal O_X)\le v(f^*g)$
for every valuation.  Thus an admissible coefficient for that divisor is
also a lower bound for the ideal threshold.

All the varieties, morphisms, rational boundaries, a rational Fano type
witness, generators, and the chosen resolutions descend to an
algebraically closed subfield $k_0$ of finite transcendence degree over
$\overline{\Q}$.  Such a field embeds in $\C$.  Normality, geometric
integrality and the identity $f_*\mathcal O_X=\mathcal O_Z$ are unchanged
by algebraically closed extension; the latter also follows from flat base
change for the proper morphism.  Projectivity and relative ampleness are
geometric properties.  Relative nefness is geometric as well: any curve
on an extension can be spread over a finite-type parameter scheme and
specialized to a curve cycle over the smaller algebraically closed field,
with the same intersection number.  Thus a negative intersection after
extension would already yield a negative intersection before extension.

For discrepancies one can use the same log resolution.  In the
nonvacuous range $0<\epsilon\le1$, a log smooth sub-boundary with all
coefficients at most $1-\epsilon$ is $\epsilon$-lc.  At a stratum this
follows by comparing with its reduced simple normal crossings support:
a primitive monomial order has at least one positive integral weight,
and each boundary weight contributes at least $\epsilon$; a
nonmonomial prime has a further nonnegative discrepancy, with an integral
contribution at least one when needed away from the strata.  This also
shows preservation of the discrepancy bound after extension.  The
coefficients and orders in the displayed resolution formula are unchanged,
so the ideal threshold is unchanged.  The complex assertion therefore
gives the same lower bound over $k_0$ and over the original field, where
the preceding general-member construction finishes the proof.
\end{proof}

\subsection{Conclusion of the two theorems}

\begin{proof}[Proof of Theorem~\ref{thm:birational}]
Over $\C$, a failure in any fixed dimension and with any fixed
$0<\epsilon\le1$ would, by Proposition~\ref{prop:H-g-input}, give the
sequence prohibited by Theorem~\ref{thm:G-pinning}.  Therefore a uniform
positive integral complement index exists.  The range $\epsilon>1$
in positive dimension is vacuous, since a prime divisor on the smooth
locus has log discrepancy one.

The birational conclusion also descends to every algebraically closed
field of characteristic zero.  Increase the complex uniform index to a
multiple greater than one, which preserves the complement assertion.
Descend a given example, its Fano type
witness, and a resolution to an algebraically closed field embeddable
in $\C$.  For the complex uniform index $n$, the complete relative
system $f_*\mathcal O_X(-nK_X)$ commutes with these field extensions.
Resolve its finitely generated fractional base ideal as in
Proposition~\ref{prop:H-saturation}.  A klt member over $\C$ implies
that every fixed-order discrepancy relevant to the fibre is strictly
positive.  Those discrepancies are unchanged by extension.  A general
member over the original infinite field consequently gives a klt
$n$-complement after shrinking, since the general residual coefficient
$1/n$ is less than one.  Its
section divisor $\Delta_s$ defines
$\Delta=\Delta_s/n\ge0$, and
$n(K_X+\Delta)=\divisor(s)$ is Cartier and linearly trivial on the
chosen neighbourhood, as required.
\end{proof}

\begin{proof}[Proof of Theorem~\ref{thm:main}]
First work over $\C$ and fix the total-space dimension $d$ and gap
$\epsilon>0$.  Theorem~\ref{thm:birational} in every dimension
$1,\ldots,d$ and for every positive discrepancy parameter, together
with its trivial dimension-zero case and the
rationalization above, verifies Chen's assertion (2) for base dimension
bound $\ell=d$.  His Theorem~1.6 therefore gives assertion (3), with a
positive constant depending only on $d$ and $\epsilon$
\cite{Chen2025}.  Every base of a contraction from a $d$-dimensional
variety has dimension at most $d$, so this applies to the given $f$.
It supplies lc on the full inverse image of a neighbourhood of $z$.
Lemma~\ref{lem:H-field-reduction}, decreasing the coefficient below one
if necessary, transports the uniform constant to an arbitrary
algebraically closed characteristic-zero field and produces a nonzero
effective Cartier divisor through $z$ there.  This proves the stated
neighbourhood assertion.
\end{proof}

\begin{proof}[Proof of Corollary~\ref{cor:real-boundaries}]
The same application of Chen's implication (2)$\Rightarrow$(3) over
$\C$ gives his full real-boundary assertion.  Only the zero-boundary
birational inputs to (2) were rationalized above.  The resulting
assertion (3) already allows arbitrary effective real $B$ with
$K_X+B$ real Cartier, and supplies the same dependence on $d$ and
$\epsilon$.  No perturbation of the nef real boundary is involved.
\end{proof}
